\documentclass[12pt, twoside]{article}
%\documentclass[12pt, twoside]{article}
\usepackage[letterpaper, margin=1in, headsep=0.2in]{geometry}
\setlength{\headheight}{0.6in}
%\usepackage[english]{babel}
\usepackage[utf8]{inputenc}
\usepackage{microtype}
\usepackage{amsmath}
\usepackage{amssymb}
%\usepackage{amsfonts}
\usepackage[nomessages]{fp} %\FPeval{\var-name}{2*sin(pi/6)}
\usepackage{siunitx} %units in math. eg 20\milli\meter
\usepackage{yhmath} % for arcs, overparenth command
\usepackage{tikz} %graphics
\usetikzlibrary{quotes, angles, arrows, arrows.meta}
\usepackage{pgfplots}
\usepgfplotslibrary{fillbetween}
\pgfplotsset{compat=1.16}
\usepackage{graphicx} %consider setting \graphicspath{{images/}}
\usepackage{parskip} %no paragraph indent
\usepackage{enumitem}
\usepackage{multicol}
\usepackage{venndiagram}

\usepackage{fancyhdr}
\pagestyle{fancy}
\fancyhf{}
\renewcommand{\headrulewidth}{0pt} % disable the underline of the header
\raggedbottom
\hfuzz=2mm %suppresses overfull box warnings

\usepackage{hyperref}

\fancyhead[LE]{\thepage}
\fancyhead[RO]{Markscheme}
\fancyhead[LO]{La Scuola d'Italia / Dr.\ Huson / IB Math 11 \\* 5 October 2026}

\begin{document}
\subsubsection*{1.8 Classwork: Compound interest and inflation --- Markscheme}

\begin{enumerate}[itemsep=0.3cm]

%--- Problem 1: Compound interest compounded annually; value, interest, first year over $8000 [6 marks] ---
\item[1.] [Maximum mark: 6]

Sofia invests \$$5000$ at $4\%$ per year compounded annually.

\medskip

\textbf{(a)} $1.04$ \hfill \textbf{A1}

\emph{Accept $1 + \dfrac{4}{100}$. A multiplier of $0.04$ or $1.4$ earns
no mark.}

\medskip

\textbf{(b)} $5000 \times 1.04^{6}$ \hfill \textbf{M1}

$=$ \$$6326.595\ldots \approx$ \$$6326.60$ \hfill \textbf{A1}

\emph{Accept \$$6330$ (3 s.f.), \$$6327$ or \$$6326.60$. Accept the TVM
solver: $N = 6$, $I\% = 4$, $PV = -5000$, $PMT = 0$, $P/Y = C/Y = 1$.}

\emph{An exponent of $5$ (giving \$$6083.26$) earns the M1 only.
Simple-interest working ($5000 + 6 \times 200 = 6200$) earns no marks.}

\medskip

\textbf{(c)} $6326.60 - 5000 = $ \$$1326.60$ \hfill \textbf{A1}

\emph{Accept \$$1330$ (3 s.f.) or \$$1327$.}

\emph{FT: award A1 for the candidate's own value from (b) minus
\$$5000$.}

\medskip

\textbf{(d)} $5000 \times 1.04^{\,n} > 8000$ \hfill \textbf{M1}

After $11$ years: \$$7697.27$ \ ($< 8000$); \quad after $12$ years:
\$$8005.16$ \ ($> 8000$)

$12$ years \hfill \textbf{A1}

\emph{Accept $1.04^{\,n} > 1.6$ solved on the GDC or with logarithms:
$n > 11.98\ldots$, so $12$ years. Accept the TVM solver for $N$ with
$PV = -5000$, $FV = 8000$.}

\emph{Accept a table search: full marks for showing both adjacent years
and stating $12$. A single row (after $12$ years the value is
\$$8005.16$, so $12$) still earns full marks, with the feedback note that
the adjacent year $11$ is what justifies ``first''. See markscheme-spec
§18.}

\emph{$11.98$ is not an acceptable final answer --- the answer must be a
whole number of years. Rounding down to $11$ earns the M1 only.}

\newpage

%--- Problem 2: Monthly compounding; real value after inflation; real value of cash [7 marks] ---
\item[2.] [Maximum mark: 7]

\$$12\,000$ invested at $5.2\%$ per year compounded monthly for $10$
years. Inflation $2.8\%$ per year.

\medskip

\textbf{(a)} $12\,000\left(1 + \dfrac{5.2}{100 \times 12}\right)^{12
\times 10}$ \hfill \textbf{M1}

$=$ \$$20\,161.668\ldots \approx$ \$$20\,161.67$ \hfill \textbf{A1}

\emph{Accept \$$20\,200$ (3 s.f.), \$$20\,162$ or \$$20\,161.67$. Accept
the TVM solver: $N = 120$, $I\% = 5.2$, $PV = -12\,000$, $PMT = 0$,
$P/Y = C/Y = 12$.}

\emph{M1 for a monthly rate of $\dfrac{5.2}{1200}$ (or $0.004333\ldots$)
and an exponent of $120$. Using $5.2\%$ compounded annually (giving
\$$19\,922.26$) earns no marks.}

\medskip

\textbf{(b)} $\dfrac{20\,161.67}{1.028^{10}}$ \hfill \textbf{M1}

$=$ \$$15\,296.61\ldots \approx$ \$$15\,296.61$ \hfill \textbf{A1}

\emph{Accept \$$15\,300$ (3 s.f.) or \$$15\,297$.}

\emph{Accept the real-interest-rate method: $5.2 - 2.8 = 2.4\%$ per year
compounded monthly, $12\,000\left(1 + \dfrac{2.4}{1200}\right)^{120} =$
\$$15\,251.33$, which also rounds to \$$15\,300$ (3 s.f.). Accept the TVM
solver with $I\% = 2.4$.}

\emph{Multiplying by $1.028^{10}$ (giving \$$26\,574$) earns no marks: the
real value must be \emph{less} than the value in (a). Subtracting
$2.8\%$ once ($0.972 \times 20\,161.67 = 19\,597$) earns no marks.}

\emph{FT: award M1A1 for the candidate's own value from (a) divided by
$1.028^{10}$.}

\medskip

\textbf{(c)} $15\,296.61 - 12\,000 = $ \$$3296.61$ \hfill \textbf{A1}

\emph{Accept \$$3300$ (3 s.f.) or \$$3297$. With the real-rate method,
\$$3251.33$ (\$$3250$ to 3 s.f.) is accepted.}

\emph{$20\,161.67 - 12\,000 = 8161.67$ earns no mark: that is the
increase before inflation is taken into account.}

\emph{FT: award A1 for the candidate's own value from (b) minus
\$$12\,000$.}

\medskip

\textbf{(d)} $\dfrac{12\,000}{1.028^{10}}$ \hfill \textbf{M1}

$=$ \$$9104.374\ldots \approx$ \$$9104.37$ \hfill \textbf{A1}

\emph{Accept \$$9100$ (3 s.f.) or \$$9104$. Accept the real-rate method:
$12\,000 \times 0.972^{10} = $ \$$9033.25$ (\$$9030$ to 3 s.f.).}

\emph{The cash keeps its face value of \$$12\,000$ but loses about
a quarter of its buying power. ``\$$12\,000$'' earns no marks.}

\end{enumerate}
\end{document}
